\input{../tex-inputs/latex-leseansicht-vorspann}

\section[Arthur Schnitzler und Paul Goldmann an Gustav Schwarzkopf, XXXX]{L04418 Arthur Schnitzler und Paul Goldmann an Gustav Schwarzkopf, XXXX}
\nopagebreak\mylabel{L04418v}
\rehead{ }\normalsize\beginnumbering\briefempfaengerindex{Schwarzkopf, Gustav@\textsc{Schwarzkopf, Gustav}!zzzGoldmann, Paul@\emph{von Paul Goldmann}!1899-09-192@{19. 9. 1899}|(be}\briefempfaengerindex{Schwarzkopf, Gustav@\textsc{Schwarzkopf, Gustav}!zzzSchnitzler, Arthur@\emph{von Arthur Schnitzler}!1899-09-192@{19. 9. 1899}|(be}
\toendnotes[C]{\smallbreak\pagebreak[2]}
\correspDesc{Versand  durch Arthur Schnitzler, Paul Goldmann am 19. 9. 1899 in Frankfurt am Main
\newline{}Übermittlung  am 19. 9. 1899 in Frankfurt am Main
\newline{}Zustellung  am 22. 9. 1899 in Wien
\newline{}Erhalt  durch Gustav Schwarzkopf am 22. 9. 1899 in Wien}\toendnotes[C]{\smallbreak}
\Standort{DLA, A:Schnitzler, HS.1985.1.1897.}
\physDesc{Bildpostkarte, 108 Zeichen
\newline{}Handschrift Arthur Schnitzler: Bleistift, deutsche Kurrent
\newline{}Handschrift Paul Goldmann: Bleistift, deutsche Kurrent
\newline{}Versand: 1) Stempel: »\nobreak{}\textcolor{gray}{1}9. 9. 99, \oindex{Frankfurt am Main@\textbf{Frankfurt am Main}|pwk}Frankfurt (Main) 1, 6–7\nobreak{}«.   2) Stempel: »\nobreak{}\oindex{I., Innere Stadt@\textbf{I., Innere Stadt}|pwk}Wien 1/1 1, 22. 9. 99, 11½V–1N., Bestellt\nobreak{}«. }\pstart{}{\pb}Herrn \textsc{Gustav
                     Schwarzkopf.}\pend{}\pstart{}Wien\oindex{Wien@\textbf{Wien}|pw}\pend{}\pstart{}\textsc{I. Tiefer Graben 23}\oindex{Wien@\textbf{Wien}!I., Innere Stadt@\textbf{I., Innere Stadt}!Tiefer Graben 23@\textbf{Tiefer Graben 23}|pw}.\pend{}{\bigskip}
\pstart
           {\pb}\textcolor{gray}{\textbf{Blick vom Deutschherrnhaus\oindex{Deutschordenskirche@\textbf{Deutschordenskirche}|pw}.}}\hfill \textcolor{gray}{\textbf{Frankfurt a. M.}}\oindex{Frankfurt am Main@\textbf{Frankfurt am Main}|pw}\pend
           \vspace{1em}
\pstart
           {\pb}Herzliche Grüße!{\\[\baselineskip]}\spacefill\mbox{ArthurSchn.}\pend
           \leftskip=0em{}\selectlanguage{ngerman}\vspace{1em}
\pstart
           \noindent{}{[}hs. Goldmann:{]} Herzlichen Gruß!\pend
           \pstart \spacefill\mbox{Dr. Paul Goldmann}\pend{}\selectlanguage{ngerman}\endnumbering\briefempfaengerindex{Schwarzkopf, Gustav@\textsc{Schwarzkopf, Gustav}!zzzGoldmann, Paul@\emph{von Paul Goldmann}!1899-09-192@{19. 9. 1899}|)be}\briefempfaengerindex{Schwarzkopf, Gustav@\textsc{Schwarzkopf, Gustav}!zzzSchnitzler, Arthur@\emph{von Arthur Schnitzler}!1899-09-192@{19. 9. 1899}|)be}\mylabel{L04418h}
\begin{anhang}
\end{anhang}\newcommand{\dateiname}{L04418}\newcommand{\titel}{Arthur Schnitzler und Paul Goldmann an Gustav Schwarzkopf, XXXX}\newcommand{\editorInnen}{Herausgegeben von Jahnke, SelmaMüller, Martin Anton}\input{../tex-inputs/latex-leseansicht-abspann}
